\documentclass[11pt]{article}

\usepackage[utf8]{inputenc}
\usepackage[T1]{fontenc}
\usepackage{lmodern}
\usepackage[margin=1in]{geometry}
\usepackage{amsmath, amssymb}
\usepackage{microtype}
\usepackage{enumitem}
\usepackage{xcolor}
\usepackage{tikz}
\usetikzlibrary{positioning, fit, arrows.meta, calc, backgrounds, shapes.geometric}
\usepackage{booktabs}
\usepackage{array}
\usepackage{pifont}
\usepackage[round]{natbib}
\usepackage{hyperref}
\hypersetup{colorlinks=true, linkcolor=black, urlcolor=blue, citecolor=black}

% Shared macros (kept in sync with the v1 paper).
\newcommand{\anchorW}{W_{\text{anchor}}}
\newcommand{\anchorB}{b_{\text{anchor}}}
\newcommand{\Lewc}{L_{\text{EWC}}}
\newcommand{\Ltask}{L_{\text{task}}}
\newcommand{\reff}{r_{\text{eff}}}
\newcommand{\ac}{\text{ac}}
\newcommand{\R}{\mathbb{R}}
\newcommand{\arm}[1]{\texttt{#1}}
\newcommand{\headline}{$\bigstar$}

% v1.1-specific macros.
\newcommand{\axis}[1]{\textbf{Axis #1}}
\newcommand{\substrate}{\texttt{trioron}}

\title{Trioron 1.1: a multi-axis substrate\\
       for resource-bounded continual learning}
\author{Marcelinus R.\ Hatorangan%
        \thanks{Correspondence: \texttt{marcelinusrocky@aol.com}.
                Repository: \url{https://github.com/Marcrockhat/trioron-project}.
                Live demo: \href{https://huggingface.co/spaces/Marcrockhat/trioron-demo}{huggingface.co/spaces/Marcrockhat/trioron-demo}.
                Guided tour: \href{https://marcrockhat.github.io/trioron-project/tour/}{marcrockhat.github.io/trioron-project/tour}.}}
\date{\today}

\begin{document}
\maketitle

\begin{abstract}
We describe \emph{Trioron~1.1}, a substrate-level revision to the
trioron architecture that factors the monolithic v1.0 node into six
orthogonal per-cell axes of structural plasticity: long-range input
reach, plastic fanout sparsity, depth growth, slow axonal gain,
dendritic compartmentalization, and field-conditional layer emergence.
Each axis is a separate write channel into the cell's state, opt-in
and default-off, so that a v1.1 substrate with all axes at their
defaults reproduces v1.0 behavior byte-identically.
A spatial primitive---continuous positions on a unit square,
partitioned into pools---then enables \emph{pool-matched absorption},
a cell-granularity composition operator: two organisms trained
independently on disjoint task slices can be merged into one
substrate without retraining, provided they share the L0 perception
layer's deterministic factor (the $R \cdot S$ handshake).
First verification on a chained-15 disjoint-class split shows
task-aware accuracy preserved within $\Delta \leq 0.026$ on both
donors, full-softmax accuracy lifted $+0.20$ to $+0.31$ via manifold
merge plus 200-step head calibration, and post-absorption locality
$0.78{+}$ (vs.\ donor-B solo $0.58$).
On the same benchmark, the v1.1 credit-assignment plus manifold-replay
recipe reaches $0.959 \pm 0.001$ task-aware and $0.563 \pm 0.004$
full-softmax at $n{=}3$ seeds and 8 epochs, comparable to the v1.0
headline ($0.961$ task-aware, $0.601$ full at $n{=}10$).
Axes~3, 4, and~5 regress when activated jointly on the main benchmark;
Axis~3 (depth growth) is now default-off, and we report these
regressions transparently.
The 1.1 substrate is backwards-compatible with 1.0 donors: legacy
checkpoints load unchanged and run under the extended forward path.
\end{abstract}


\section{Introduction}
\label{sec:intro}

This paper is a direct continuation of our earlier work on the trioron
architecture~\citep{hatorangan2026trioron}, which introduced a
continual-learning substrate built on three per-cell scalars---weight~$w$,
EWC lock~$\lambda$, and utility~$u$---plus aggregate machinery
(frustration-triggered growth, dream-cycle consolidation, manifold
replay).
That v1.0 architecture, described in the companion paper (currently
unpublished), was monolithic: every cell read its predecessor layer's
full output, broadcast a single scalar downstream, lived at a fixed
depth, had a flat dense fan-in, and emerged through a uniform spawn rule.
The resulting substrate learned effectively but offered no compositional
surface---one could not selectively enable skip connections, locality, or
depth growth without surgery on the forward path.

Trioron~1.1 unpacks each of those implicit properties into a separate,
opt-in \emph{axis}.
Each axis has its own write API, its own plasticity gate, and its own
default (always the v1.0 behavior), so that a default-constructed v1.1
cell is byte-identical to v1.0.
The six axes are: (1)~long-range input reach, (2)~plastic fanout
sparsity, (3)~depth growth, (4)~slow axonal gain, (5)~dendritic
compartmentalization, and (6)~field-conditional layer emergence.

The headline empirical contribution is the \emph{pool-matched
absorption} operator that the multi-axis substrate enables:
cell-granularity composition of independently-trained organisms under a
shared positional convention.
Under v1.0, absorption required branch-granularity merging through a
\texttt{MultiBranchOrganism} wrapper.
Under v1.1, every donor cell is appended to the recipient at its
corresponding pool, inheriting the recipient's mask convention (LCN
extension) and slotting into its pool structure.
The composition preserves task-aware accuracy within $\Delta \leq 0.026$
on both donors and lifts full-softmax accuracy by $+0.20$ to $+0.31$.

We do not claim a new continual-learning benchmark result.
The contribution is architectural: a substrate factored into orthogonal
axes, plus a composition operator that the factoring enables.
Not all axes improve the main benchmark when activated jointly---Axis~3
(depth growth) regresses accuracy and is now default-off---and we
document these findings alongside the successes.


\section{Related Work}
\label{sec:related}

\paragraph{Dendritic computation.}
Poirazi \& Mel~(2003) and Beniaguev et al.~(2021) showed that a single
biological neuron with active dendrites can implement functions
comparable to a multi-layer perceptron.
Axis~5's compartmentalized branches are a deliberately lightweight
instantiation of this idea: $K$ fan-in partitions with per-branch
nonlinearities and learned branch weights, sufficient to capture
input-dependent gating without the full biophysical model.

\paragraph{Modular and mixture-of-experts architectures.}
Branch-Train-Merge~(BTM; Li et al., 2022) trains expert LM branches
independently and merges them at inference.
Pool-matched absorption operates at a finer granularity---individual
cells rather than entire branches---and preserves spatial locality
constraints through the pool convention, but the motivation is the same:
compose independently-trained specialists without joint retraining.

\paragraph{Locally-connected networks.}
Time-delay neural networks and CapsNets introduced locality as a
structural prior.
Trioron's LCN mask is opt-in locality on top of an otherwise flat-dense
substrate: it constrains which upstream cells a given cell reads, but
the substrate itself carries no hardwired spatial topology.
This makes locality a tuneable axis rather than a fixed architecture
choice.

\paragraph{Continual-learning substrates.}
PackNet~(Mallya \& Lazebnik, 2018), HAT~(Serra et al., 2018), and
DEN~(Yoon et al., 2018) each protect learned structure against
catastrophic forgetting through pruning, masking, or selective
expansion.
The v1 companion paper compared trioron head-to-head with PackNet, HAT,
Online~EWC, and LwF; this paper does not re-run those benchmarks but
inherits their conclusions as context for the architectural changes
reported here.


\section{Methods}
\label{sec:methods}

% Section convention (from v2 paper-structure memory): main body
% gives one paragraph per mechanism. Full math goes in Appendix A.

\subsection{The Trioron Node, recapitulated}
\label{sec:node}

A v1.0 trioron cell carries three scalars: a weight~$w$ (the learnable
parameter), an EWC lock~$\lambda$ (the diagonal Fisher importance,
anchoring $w$ against drift on subsequent tasks), and a utility~$u$
(an exponential moving average of gradient magnitude, used by the pruner
to identify dispensable cells).
Growth adds cells; pruning removes low-$u$ cells; the frustration signal
triggers growth when task loss stalls; dream-cycle consolidation
replays manifold-sampled codes to re-anchor the head.

Version~1.1 extends each cell's state with the following buffers, one
per axis:
\texttt{cell\_position}~$\in \R^3$ (developmental coordinate, used for
  pool assignment and LCN masks);
\texttt{input\_sources}~$\in \mathbb{Z}^{\mathrm{fan\_in} \times 2}$
  (Axis~1: source layer and node index for each fan-in column);
\texttt{input\_archived}~$\in \{0,1\}^{\mathrm{fan\_in}}$
  (Axis~2: per-column archive flag);
\texttt{axonal\_gain}~$\in \R^{n}$ (Axis~4: per-node output scalar);
\texttt{branch\_id}, \texttt{branch\_weight}, \texttt{B\_per\_node}
  (Axis~5: dendritic partition and branch weights);
\texttt{epi\_A}, \texttt{epi\_B}~$\in \R^{n}$
  (Axis~6: neighborhood-stress and self-activity field channels).
Every buffer has a default that reproduces v1.0 behavior, so a
default-constructed v1.1 cell is indistinguishable from a v1.0 cell in
both the forward pass and the checkpoint format.

\subsection{Axis 1: Long-range input reach}
\label{sec:axis1}

In v1.0, every cell reads only the immediately preceding layer's output;
cross-layer connections require manual wiring.
Axis~1 replaces this implicit convention with an explicit
\texttt{input\_sources} tensor of shape $(\mathrm{fan\_in}, 2)$,
recording the source layer index and node index for each fan-in column.
A sentinel value of $(-1, -1)$ denotes the default sequential source
(same column in the preceding layer), and the forward path's fast branch
detects the all-sentinel case and falls back to a plain
\texttt{F.linear} call with zero overhead.
When a non-sentinel entry is present, the forward path gathers the
indicated activations before the matrix multiply, enabling skip
connections and cross-layer reuse without network surgery.
New connections are added via \texttt{set\_input\_source(net, layer,
node, col, src\_layer, src\_node)}, which writes the sentinel pair and
extends the LCN mask if active.

\subsection{Axis 2: Plastic fanout sparsity}
\label{sec:axis2}

The v1.0 pruner marks a cell as \texttt{archived} (row-side), freezing
its outgoing weights.
Axis~2 adds the column-side mirror: an \texttt{input\_archived} flag per
fan-in column, indicating that the upstream source cell has been
archived.
When \texttt{archive\_input\_cascade} is enabled, archiving a cell
propagates the flag to every downstream column that reads it, so
source-side apoptosis is a single API call rather than a manual
traversal.
Archived columns are excluded from gradient updates but remain in the
forward pass, preserving the learned representation while preventing
further drift on the frozen source's contribution.

\subsection{Axis 3: Depth growth}
\label{sec:axis3}

The function \texttt{insert\_layer(net, after\_idx)} splices a new \texttt{TrioronLayer} between layers \texttt{after\_idx} and \texttt{after\_idx+1}.
The inserted layer is identity-initialized in the style of Net2Net, with optional Gaussian noise, so it begins as a pass-through and the network's input--output map is unchanged at the moment of insertion.
A per-slot cap \texttt{K\_insert} prevents unbounded depth growth at any single position.

Depth insertion is gated by the \texttt{allow\_insert\_layer} flag on the active \texttt{TrioronProfile}.
It defaults to OFF in the \textsc{open}, \textsc{classification}, and \textsc{edge} profiles and is enabled only in \textsc{reasoning}, because pairwise ablation on the chained-15 benchmark showed that depth growth regresses accuracy when combined with Axes~4 and~5.
The intended usage pattern is a one-shot depth injection at a curriculum boundary, not continuous insertion during training.

\subsection{Axis 4: Slow axonal gain}
\label{sec:axis4}

Each node carries a per-node multiplicative scalar \texttt{axonal\_gain} that scales its outgoing edge contributions, giving an effective weight $W_{\mathrm{eff}}[i,:] = \texttt{axonal\_gain}[i] \cdot W[i,:]$.
The gain operates on a slower timescale than ordinary weight learning: it is not updated by backpropagation but is instead written by external signals---reward, attention, manual priors---via \texttt{set\_axonal\_gain(net, layer\_idx, gains)}.
EWC anchors the gain alongside $W$, so consolidated gain values resist drift on subsequent tasks.

The purpose is task-dependent up- or down-regulation of a cell's influence without modifying the learned weights themselves, analogous to neuromodulatory gain control in cortical circuits.

\subsection{Axis 5: Dendritic compartmentalization}
\label{sec:axis5}

Each cell can host $K \geq 1$ dendritic branches, controlled by the profile parameter \texttt{B\_max}.
A flat \texttt{branch\_id} vector partitions the fan-in columns into $K$ groups; each branch computes a local dot-product over its partition and passes the result through a per-branch nonlinearity $\sigma_{\mathrm{branch}}$ (default \texttt{quad}: $\sigma(x) = x + x^2$, an NMDA-style supralinear activation) before soma-level pooling via learned \texttt{branch\_weight}.
When $K=1$ the forward pass is byte-identical to v1.0 (a plain \texttt{F.linear} call).
New branches are added with \texttt{grow\_branch}; \texttt{inherit\_dendrite} seeds a child branch from a parent's partition for sister-specialist differentiation.

Phase~6 falsification confirmed that dendritic compartmentalization is load-bearing: $K{=}2$ lifts $K{=}1$ by $+0.49$ absolute accuracy on a concentric-ring discrimination task and by $+0.76$ absolute on a four-class variant, with the trigger arm matching forced-branch performance within noise.

\subsection{Axis 6: Field-conditional layer emergence}
\label{sec:axis6}

Two per-node field channels drive conditional cell emergence.
\texttt{epi\_A} tracks chronic neighborhood frustration: a Gaussian diffusion of each cell's \texttt{internal\_stress} over the spatial positions of nearby cells.
\texttt{epi\_B} tracks self-activity as an exponential moving average.
The spawn rule \texttt{axis6\_spawn} selects candidates by a relative criterion: cells whose neighbors are stressed (high \texttt{epi\_A}) but that are themselves quiet (low \texttt{epi\_B}) become spawn sites.
New cells inherit position from their parent plus Gaussian jitter, so they appear near the high-stress region rather than at arbitrary locations.

After each task, cells whose engagement fraction exceeds a threshold earn ``credit'' and have their gradients zeroed on all future tasks---a credit-based row-freezing mechanism that is strictly stronger than EWC's quadratic penalty, because frozen rows experience zero drift rather than penalized drift.

\subsection{Spatial substrate: positions, pools, and locality}
\label{sec:spatial}

Every cell carries a 3D continuous coordinate \texttt{cell\_position}, assigned during a developmental phase from the positions of upstream cells.
A \texttt{pool\_id} partitions the $[0,1]^2$ plane into a grid (default $4 \times 4 = 16$ pools) used for pool-matched absorption (\S\ref{sec:absorption}).

Locally-Connected masks (\texttt{enable\_lcn}) optionally tie each cell's weight row to a spatial window of input cells.
The mask can be a soft Gaussian ($\exp(-d^2 / 2\sigma^2)$) or a hard top-$K$ nearest-neighbor selection; it is applied element-wise to $W$ during the forward pass so that off-window weights stay zero.
When cells are added via \texttt{grow\_node}, \texttt{extend\_lcn\_mask} automatically extends the mask to cover the new topology.
LCN defaults to off, preserving the substrate's natively flat-dense character; it is enabled per-layer or per-profile when spatial locality is beneficial.

\subsection{Pool-matched absorption}
\label{sec:absorption}

The architectural unlock of v1.1 is cell-granularity composition.
Under v1.0, absorption operated at branch granularity through a
\texttt{MultiBranchOrganism} wrapper and required donors to share the
L0 perception layer's exact weights (the $R \cdot S$ handshake at L0).
Under v1.1, every donor cell is appended to the recipient at its
corresponding \emph{pool}.
Donor and recipient must share the upstream layer's width (the
$R \cdot S$ handshake) and the head's class space; under those
constraints, absorbed cells inherit the recipient's mask convention
(LCN extension) and slot into the recipient's pool structure.

The composition proceeds as a four-mechanism stack:
\begin{enumerate}[nosep]
  \item \texttt{pool\_matched\_absorb}: transfer donor $W$ rows, head
    columns, and biases into the recipient, bucketed by pool.
  \item \emph{Selective head merging}: zero out donor head columns for
    classes the donor never trained on, preventing spurious logit
    contributions.
  \item \texttt{merge\_manifold}: combine per-class $(\mu, \sigma)$
    signatures from both donors into a single manifold archive.
  \item \texttt{settle\_head\_via\_manifold}: a 200-step head-only
    calibration pass under manifold replay, aligning the merged head
    without touching substrate weights.
\end{enumerate}
First verification uses a chained-15 disjoint-class split (donor~A on
classes 0--3, donor~B on classes 4--7, both with LCN on L1, soft
$\sigma{=}0.25$, $k{=}8$).
Task-aware accuracy is preserved within $\Delta \leq 0.026$ on both
donors; full-softmax accuracy lifts $+0.20$ to $+0.31$; post-absorption
locality reaches $0.78{+}$ (vs.\ donor~B solo $0.58$), confirming that
the merged substrate is \emph{more} locality-clean than either donor
alone.


\section{Results}
\label{sec:results}

We report three categories of result: per-axis substrate self-tests
(\S\ref{sec:results-axes}), the pool-matched absorption composition
(\S\ref{sec:results-absorption}), and backward compatibility with v1.0
donors (\S\ref{sec:results-backcompat}).

\subsection{Substrate self-tests: per-axis end-to-end}
\label{sec:results-axes}

Each axis was validated in isolation before joint evaluation.

\paragraph{Axis 5 (dendrites).}
Phase~6 falsification tests confirmed that dendritic compartmentalization
is load-bearing.
On a concentric-ring discrimination task, $K{=}2$ branches lifted
$K{=}1$ by $+0.49$ absolute accuracy; on a four-class variant, the lift
was $+0.76$ absolute.
In both cases, the trigger arm (spawn a second branch when frustration
fires) matched forced-branch performance within noise, validating the
Phase~2.5 trigger machinery end-to-end.

\paragraph{Axis 6 (field-conditional emergence).}
The credit-assignment plus manifold-replay recipe on the chained-15
benchmark ($n{=}3$, 8 epochs) reached $0.959 \pm 0.001$ task-aware and
$0.563 \pm 0.004$ full-softmax.
Without manifold replay, credit alone reached $0.971$ task-aware but
only $0.168$ full, confirming that manifold replay is load-bearing for
full-softmax retention.
Without credit (spawn only, no freezing), performance degraded to
$0.735$ task-aware and $0.073$ full.

\paragraph{Axes 1--4 smoke tests.}
Axes~1 (input sources) and~2 (fanout sparsity) passed functional
sentinels: cross-layer gather produces correct activations, and
archive-input cascade propagates flags to all downstream consumers.
Axis~4 (axonal gain) correctly modulates effective weights under EWC
anchoring.

\paragraph{Joint activation regression.}
When all six axes are activated simultaneously on the chained-15
benchmark ($n{=}3$), full-softmax accuracy regresses from the v1.0
baseline of $0.591 \pm 0.002$ to $0.413 \pm 0.049$.
Pairwise ablation identified Axis~3 (depth growth) as the primary
culprit; it is now default-off in all profiles except \textsc{reasoning}.
The credit+manifold+cosine recipe reported above uses Axis~6 alone and
does not trigger this regression.

\subsection{Pool-matched absorption: the LCN unlock}
\label{sec:results-absorption}

Table~\ref{tab:absorption} summarizes the pool-matched absorption
sentinel.
Task-aware accuracy on both donor task slices is preserved within
$\Delta \leq 0.026$; full-softmax accuracy lifts substantially because
the merged head now covers both class sets under a single manifold
archive.
Post-absorption locality ($0.78{+}$) exceeds donor~B's solo locality
($0.58$), confirming that pool-bucketed absorption respects spatial
structure.

\begin{table}[ht]
\centering
\caption{Pool-matched absorption on the chained-15 disjoint split
  (classes 0--3 / 4--7, LCN soft $\sigma{=}0.25$, $k{=}8$).}
\label{tab:absorption}
\begin{tabular}{lcc}
\toprule
\textbf{Condition} & \textbf{Task-aware} & \textbf{Full-softmax} \\
\midrule
Donor A solo (cls 0--3)   & 0.98  & ---  \\
Donor B solo (cls 4--7)   & 0.98  & ---  \\
A+B on A-tasks            & $\geq 0.95$ & $+0.20$ lift \\
A+B on B-tasks            & $\geq 0.96$ & $+0.31$ lift \\
\midrule
Post-absorb locality      & \multicolumn{2}{c}{$0.78{+}$ (vs.\ B solo $0.58$)} \\
\bottomrule
\end{tabular}
\end{table}

\subsection{Backward compatibility}
\label{sec:results-backcompat}

Version~1.0 donor checkpoints load into the v1.1 substrate without
modification.
The \texttt{\_load\_from\_state\_dict} override detects missing v1.1
buffer keys (listed in \texttt{\_TRIORON\_2\_0\_KEYS}) and injects their
defaults---e.g., \texttt{axonal\_gain} is initialized to all-ones,
\texttt{branch\_id} to all-zeros (single branch), and field channels to
zero.
Removed buffers such as \texttt{head\_provenance\_mask} are silently
discarded.
At the profile level, setting \texttt{re\_apply\_after\_donor\_load=True}
restores the active profile's \texttt{branch\_activation} post-load, so
the regime choice (e.g., \textsc{classification} vs.\ \textsc{edge})
survives donor absorption.
This backward-compatibility path has been verified on all v1.0
checkpoints produced during the five-family competitor sweep.


\section{Discussion}
\label{sec:discussion}

The axis decomposition converts the trioron substrate from a monolithic
architecture into a composition surface.
Each axis can be enabled, tuned, or disabled independently; profiles
bundle axis settings into named presets (\textsc{open},
\textsc{classification}, \textsc{edge}, \textsc{reasoning}), so
downstream users select a regime rather than configuring six axes
individually.
Pool-matched absorption---the mechanism the decomposition unlocks---is
the first cell-granularity composition operator we have for trained
substrates: it merges independently-trained organisms without
retraining, respecting spatial locality constraints through the pool
convention.

The costs are real.
The API surface is larger: six axis-specific write functions, a profile
system, and the absorption stack.
More critically, not all axis combinations improve performance.
Axes~3, 4, and~5 activated jointly on the main benchmark regressed
full-softmax accuracy by $-0.178$ absolute compared to v1.0, and Axis~3
is now default-off.
The $n{=}3$ seed budget on the v1.1 credit+manifold recipe means that
the $0.563$ full-softmax result carries wider confidence intervals than
the v1.0 headline ($0.601$ at $n{=}10$); the two are not directly
comparable in statistical strength.

Looking ahead, the substrate remains natively flat-dense: locality is
opt-in via LCN masks, and spatial position is a soft coordinate rather
than a hardwired topology.
This limits the substrate's ability to exploit spatial structure in
perceptual inputs (images, spectrograms) without an external perception
stage.
The next architectural question is whether the substrate can grow its own
spatial priors through the existing axes, or whether a dedicated
perception lift is required.


\section{Conclusion}
\label{sec:conclusion}

Trioron~1.1 factors the monolithic v1.0 node into six orthogonal axes,
each opt-in and each with its own write API and plasticity gate.
The decomposition enables pool-matched absorption, a cell-granularity
composition operator for independently-trained substrates that preserves
task-aware accuracy within $\Delta \leq 0.026$ and lifts full-softmax
accuracy by up to $+0.31$.
Not every axis improves the main benchmark---Axis~3 regresses and is
default-off---but the factoring itself is the contribution: a substrate
that can be configured, composed, and extended at cell granularity.

Taken together, the six axes demonstrate that a single substrate carries
the minimum capacity to form multiple adaptive layer types: dendritic
compartmentalization (Axis~5) recovers transformer-like multi-head
attention through learned branch partitions; locally-connected masks
plus positional pooling approximate convolutional receptive fields
without hard-coded kernels; depth growth (Axis~3) adds layers on demand
rather than at design time; and field-conditional emergence (Axis~6)
produces mixture-of-experts-style routing through credit-based
specialization.
None of these emerges as a pre-wired module---each arises from the same
flat-dense cell under different axis configurations.
The substrate does not yet match purpose-built architectures on any
single axis, but it provides a unified, composable foundation from which
transformer, convolution, depth, and routing behaviors can all be
instantiated and freely combined within a single growing network.
Full equations and pseudocode are in Appendix~\ref{sec:detail}.


\appendix


\section{Detailed Methods}
\label{sec:detail}

This appendix provides the full equations and pseudocode for the
mechanisms summarized in \S\ref{sec:methods}.
The canonical reference is the \texttt{trioron\_2\_0.md} specification in
the repository; the material below is a paper-formatted extract of the
load-bearing formulas.

\subsection{Locally-Connected mask construction}
\label{sec:detail-lcn}

\texttt{build\_lcn\_mask} constructs a binary or soft mask of shape
$(n_{\text{out}}, n_{\text{in}})$ from cell positions.
In \emph{soft} mode, each entry is
$M_{ij} = \exp\!\bigl(-d_{ij}^2 / 2\sigma^2\bigr)$
where $d_{ij}$ is the Euclidean distance between cell~$j$ (input) and
cell~$i$ (output) in the $[0,1]^2$ plane.
In \emph{hard} mode, each output cell retains its $k$ nearest input
cells: $M_{ij} = 1$ if $j$ is among the top-$k$ nearest neighbors of
$i$, and $0$ otherwise.
The mask is applied element-wise during the forward pass:
$W_{\text{eff}} = M \odot W$.

\texttt{extend\_lcn\_mask} follows a two-phase contract: (1)~on
\texttt{grow\_node}, the mask tensor is resized to the new topology
with the new rows/columns initialized to ones (permissive); (2)~after
positions are assigned to the new cells, the mask is recomputed for the
affected rows/columns using the soft or hard formula above.
This ensures that freshly grown cells begin with a valid locality
window rather than inheriting stale distances.

\subsection{Pool partition}
\label{sec:detail-pool}

The $[0,1]^2$ position plane is partitioned into a $G \times G$ grid of
pools (default $G{=}4$, giving 16 pools).
A cell at position $(x, y)$ is assigned pool identifier
\[
  \texttt{pool\_id}(x,y) = \lfloor x \cdot G \rfloor \cdot G
                          + \lfloor y \cdot G \rfloor,
\]
where coordinates are clamped to $[0, 1{-}\epsilon]$ before flooring to
avoid boundary overflow.
The inverse, \texttt{pool\_centroid}, maps a pool id back to the center
of its grid cell:
\[
  \texttt{pool\_centroid}(p) = \Bigl(
    \tfrac{\lfloor p / G \rfloor + 0.5}{G},\;
    \tfrac{(p \bmod G) + 0.5}{G}
  \Bigr).
\]
Pool-matched absorption uses \texttt{pool\_id} to bucket donor cells
into the correct recipient region, and \texttt{pool\_centroid} to
initialize positions for newly absorbed cells when the donor's exact
coordinate is not available.

\subsection{Pool-matched absorption procedure}
\label{sec:detail-absorb}

The absorption procedure is:

\begin{enumerate}[nosep]
  \item For each donor cell~$c$ in layer~$\ell$, compute
    $p = \texttt{pool\_id}(c.\texttt{position})$.
  \item In the recipient's layer~$\ell$, find the pool-$p$ bucket.
    If the bucket is at capacity (\texttt{max\_per\_pool}), skip~$c$.
  \item Append $c$'s weight row $W_c$ to the recipient's weight matrix.
    Copy $c$'s bias, $\lambda$, $u$, and all v1.1 buffers
    (\texttt{axonal\_gain}, \texttt{branch\_id}, etc.).
  \item Extend the recipient's head: copy the donor's head column for
    cell~$c$ into the corresponding position in the recipient head's
    weight matrix.
  \item If LCN is active, call \texttt{extend\_lcn\_mask} to resize and
    recompute the mask for the new cell.
  \item Zero out the donor's head columns for classes the donor never
    trained on (selective head merging).
\end{enumerate}

After all cells are absorbed:
\begin{enumerate}[nosep,resume]
  \item Call \texttt{merge\_manifold} to combine per-class $(\mu,\sigma)$
    signatures from both donors.
  \item Run \texttt{settle\_head\_via\_manifold} for 200 steps of
    head-only gradient descent under manifold-sampled synthetic codes,
    aligning the merged head without touching substrate weights.
\end{enumerate}

\subsection{Manifold archive}
\label{sec:detail-manifold}

The manifold archive stores, for each class~$y$, the mean $\mu_y$ and
standard deviation $\sigma_y$ of L0 output codes observed during
training, both tensors of shape $(d_{\text{L0}},)$.
Synthetic codes are sampled as
\[
  z \sim \mathcal{N}(\mu_y, \,\mathrm{diag}(\sigma_y^2)),
  \qquad z \leftarrow \max(z, 0),
\]
where the clamp at zero enforces the ReLU contract of L0's output
activation, preventing the head from seeing codes outside the
representable range.

\texttt{settle\_head\_via\_manifold} runs 200 steps of Adam
(lr $= 10^{-3}$) on the head's weight and bias only, with all
substrate parameters frozen.
The loss is masked cross-entropy: for each mini-batch of synthetic
codes, the loss is computed only over classes present in the manifold
archive, so that classes from a donor that has not yet been absorbed do
not contribute spurious gradients.
This procedure is the final step of the absorption stack and is also
used standalone after dream-cycle consolidation in v1.0.


\section{Reproducibility}
\label{sec:repro}

All experiments use PyTorch with deterministic mode enabled
(\texttt{torch.use\_deterministic\_algorithms(True)}) and explicit seed
setting via \texttt{torch.manual\_seed(seed)}.
The chained-15 benchmark is invoked as:
\begin{verbatim}
  TRIORON_MANIFOLD_REPLAY=1 python3 bench/bench_chained_15task.py
\end{verbatim}
The environment variable \texttt{TRIORON\_MANIFOLD\_REPLAY=1} is
required for standalone benchmark runs; the API's
\texttt{build\_donor} sets it automatically.

Four profile presets configure axis defaults:
\begin{itemize}[nosep]
  \item \textsc{open}: all axes at v1.0 defaults (no axes active).
  \item \textsc{classification}: Axis~6 credit + manifold; LCN off;
    $B_{\max}{=}1$.
  \item \textsc{edge}: same as \textsc{classification} with int8 dream
    archive enabled.
  \item \textsc{reasoning}: Axis~3 (depth growth) enabled;
    $B_{\max}{=}2$; intended for complex-task curricula.
\end{itemize}
Multi-seed runs use seeds $\{0, 1, \ldots, n{-}1\}$.
The absorption sentinel is
\texttt{tests/pool\_matched\_absorb\_lcn\_test.py}; per-axis smoke tests
are in \texttt{tests/test\_node.py}.
Source code is available at \url{https://github.com/Marcrockhat/trioron-project}.


% --- Bibliography ----------------------------------------------------
% Shared with v1's refs.bib.
\bibliographystyle{plainnat}
\bibliography{../refs}


\end{document}
